% Part: sets-functions-relations
% Chapter: relations
% Section: relations-as-sets

\documentclass[../../../include/open-logic-section]{subfiles}

\begin{document}

\olfileid{sfr}{rel}{set}
\olsection{సమితులుగా సంబంధాలు}

\begin{explain}
\olref[sfr][set][imp]{sec}లో $\Nat$, $\Int$, $\Rat$, $\Real$ అనే
కొన్ని ముఖ్యమైన సమితులను పేర్కొన్నాం. ఈ సమితులలో కొన్నింటి
\tecase{gen}{!!{element}s} మధ్య ఆసక్తికరమైన సంబంధాలు మీకు తప్పక గుర్తుంటాయి.
ఉదాహరణకు వీటిలో ప్రతి సమితిపైనా ఒక పూర్తిగా ప్రామాణికమైన
\emph{క్రమ సంబంధం} ఉంది. అలాగే ప్రతి వస్తువుకూ తనతోనే ఉండి, మరే ఇతర
వస్తువుతోనూ ఉండని \emph{తాదాత్మ్య సంబంధం} ఉంది. మనం మరెన్నో ఆసక్తికరమైన
సంబంధాలను చూస్తాం; సాధ్యమయ్యే సంబంధాలు ఇంకా ఎక్కువే ఉన్నాయి.
వాటిని పరిశీలించే ముందు, సంబంధాలను ఒక ప్రత్యేకమైన రకం సమితులుగా
చూడవచ్చని వివరిద్దాం.
  
దీనికోసం \olref[sfr][set][pai]{sec}లోని రెండు విషయాలను గుర్తుచేసుకుందాం.
మొదటిది \emph{క్రమిత జత}: $a$, $b$ ఇచ్చినప్పుడు $\tuple{a, b}$ను
ఏర్పరచవచ్చు. ఇక్కడ మూలకాల క్రమం \emph{ముఖ్యమే}. కనుక $a \neq b$ అయితే
$\tuple{a, b} \neq \tuple{b, a}$. (దీనిని క్రమరహిత జతలతో, అంటే
$2$-మూలకాల సమితులతో పోల్చండి: వాటికి $\{a, b\}=\{b, a\}$.)
రెండవది \emph{కార్టీజియన్ లబ్ధం}: $A$, $B$ సమితులు అయితే,
$x \in A$, $y \in B$ అయ్యే $\tuple{x, y}$ జతలన్నింటి సమితి
$A \times B$ను ఏర్పరచవచ్చు. ప్రత్యేకించి, $A^{2}= A \times A$
అనేది $A$లోని మూలకాలతో ఏర్పడే క్రమిత జతలన్నింటి సమితి.
  
ఇప్పుడు ఒక సమితిపై ఒక నిర్దిష్ట సంబంధాన్ని పరిశీలిద్దాం: సహజ సంఖ్యల
సమితి $\Nat$పై $<$-సంబంధం. $n<m$ అయ్యే సంఖ్యల జతలు
$\tuple{n, m}$ అన్నింటి సమితిని తీసుకుందాం. అంటే,
\[
  R=\Setabs{\tuple{n, m}}{n, m \in \Nat \text{ మరియు } n<m}.
\]
$n$ అనేది $m$ కంటే చిన్నదిగా ఉండటానికీ, $\tuple{n, m}$ అనే జత
$R$లో సభ్యంగా ఉండటానికీ దగ్గరి సంబంధం ఉంది. అది:
\[
      n<m\text{ అయితే, అయితేనే }\tuple{n, m} \in R.
\]
నిజానికి ఏ సమాచారమూ పోకుండా, $R$ అనే సమితినే $\Nat$పై
$<$-సంబంధంగా \emph{పరిగణించవచ్చు}.

ఇలాగే సంఖ్యల మధ్య ఏ సంబంధానికైనా $\Nat^{2}$లో ఒక ఉపసమితిని
నిర్మించవచ్చు. తిరుగుగా, సంఖ్యల జతల సమితి $S \subseteq \Nat^{2}$
ఏది ఇచ్చినా, దానికి అనుగుణమైన సంబంధం సంఖ్యల మధ్య ఉంటుంది:
$\tuple{n, m} \in S$ అయితే, అయితేనే $n$కూ $m$కూ ఆ సంబంధం
ఉండేలా తీసుకోవాలి. దీనివల్ల ఈ నిర్వచనం సమర్థించబడుతుంది:
\end{explain}
  
\begin{defn}[ద్విస్థానిక సంబంధం] 
$A$ అనే సమితిపై ఒక \emph{ద్విస్థానిక సంబంధం} అంటే $A^{2}$లో ఒక
ఉపసమితి. $R \subseteq A^{2}$ అనేది $A$పై ద్విస్థానిక సంబంధమై,
$x, y \in A$ అయితే, $\tuple{x, y} \in R$కు బదులుగా కొన్నిసార్లు
$Rxy$ (లేదా $xRy$) అని రాస్తాం.
\end{defn}
  
\begin{ex}
  \ollabel{relations}
సహజ సంఖ్యల జతల సమితి $\Nat^{2}$ను ఈ విధంగా రెండు దిశలలో
మాత్రిక రూపంలో జాబితా చేయవచ్చు:
\[
  \begin{array}{ccccc}
  \mathbf{\tuple{ 0,0 }} & \tuple{ 0,1 } &
    \tuple{ 0,2 } & \tuple{ 0,3 } & \ldots\\
  \tuple{ 1,0 } & \mathbf{\tuple{ 1,1 }} &
    \tuple{ 1,2 } & \tuple{ 1,3 } & \ldots\\
  \tuple{ 2,0 } & \tuple{ 2,1 } &
    \mathbf{\tuple{ 2,2 }} & \tuple{ 2,3 } & \ldots\\
  \tuple{ 3,0 } & \tuple{ 3,1 } & \tuple{ 3,2 } &
    \mathbf{\tuple{ 3,3 }} & \ldots\\
  \vdots & \vdots & \vdots & \vdots & \mathbf{\ddots}
  \end{array}
\]
ఇక్కడ వికర్ణాన్ని ముదురు అక్షరాలతో చూపాం. ఎందుకంటే వికర్ణంపై ఉన్న
జతలతో ఏర్పడే $\Nat^2$ ఉపసమితి, అంటే
\[
  \{\tuple{0,0 }, \tuple{ 1,1 }, \tuple{ 2,2 }, \dots\},
  \]
అనేది $\Nat$పై \emph{తాదాత్మ్య సంబంధం}. (తాదాత్మ్య సంబంధాన్ని
తరచుగా ఉపయోగిస్తాం కనుక, ఏ సమితి $A$కైనా
$\Id{A}=\Setabs{\tuple{ x,x }}{x \in A}$ అని నిర్వచిద్దాం.)
వికర్ణానికి పైన ఉన్న జతలన్నింటి ఉపసమితి, అంటే
\[
  L = \{\tuple{ 0,1 },\tuple{ 0,2 },\ldots,\tuple{ 1,2 },
  \tuple{ 1,3 }, \dots, \tuple{ 2,3 }, \tuple{ 2,4 },\ldots\},
\]
అనేది \emph{కంటే చిన్నది} అనే సంబంధం. అంటే $Lnm$ కావడానికి
అవసరమైన మరియు సరిపడే షరతు $n<m$ అవడం. వికర్ణానికి కింద ఉన్న
జతల ఉపసమితి, అంటే
\[
  G=\{ \tuple{ 1,0 },\tuple{ 2,0 },\tuple{
    2,1 }, \tuple{ 3,0 },\tuple{ 3,1 },\tuple{ 3,2 }, \dots\},
\]
అనేది \emph{కంటే పెద్దది} అనే సంబంధం. అంటే $Gnm$ కావడానికి
అవసరమైన మరియు సరిపడే షరతు $n>m$ అవడం. $L$కూ $I$కూ సమ్మేళనమైన
$K=L\cup I$ అనేది \emph{చిన్నది లేదా సమానమైనది} అనే సంబంధం:
$Knm$ కావడానికి అవసరమైన మరియు సరిపడే షరతు $n \le m$ అవడం.
అలాగే, $H=G \cup I$ అనేది \emph{పెద్దది లేదా సమానమైనది} అనే సంబంధం.
ఈ $L$, $G$, $K$, $H$ సంబంధాలు \emph{క్రమాలు} అనే ప్రత్యేక రకానికి
చెందుతాయి. $L$, $G$లకు ఈ ధర్మం ఉంది:
ఏ సంఖ్యకూ తనతో $L$ గానీ $G$ గానీ ఉండదు
(అంటే ప్రతి $n$కూ $Lnn$, $Gnn$ రెండూ అసత్యం).
ఈ ధర్మం ఉన్న సంబంధాలను \emph{అస్వావర్తన} సంబంధాలు అంటాం.
అవి క్రమాలు కూడా అయితే \emph{కఠిన క్రమాలు} అంటాం.
\end{ex}

\begin{explain}
క్రమాలు, తాదాత్మ్యం ముఖ్యమైన సహజ సంబంధాలే అయినా, మన నిర్వచనం ప్రకారం
$A^{2}$లోని \emph{ఏ} ఉపసమితి అయినా $A$పై సంబంధమే అని గుర్తుంచుకోవాలి;
అది ఎంత అసహజంగా లేదా కృత్రిమంగా కనిపించినా సరే.
ప్రత్యేకించి, ఏ సమితిపైనైనా $\emptyset$ ఒక సంబంధం
(\emph{శూన్య సంబంధం}: ఏ మూలకాల జతకూ అది ఉండదు).
$A^{2}$ కూడా $A$పై ఒక సంబంధమే; ప్రతి జతకూ ఉండే ఈ సంబంధాన్ని
\emph{సార్వత్రిక సంబంధం} అంటాం.
అలాగే $E=\Setabs{\tuple{n, m}}{n>5 \text{ లేదా } m \times n \ge 34}$
వంటిది కూడా సంబంధంగానే లెక్కలోకి వస్తుంది.
\end{explain}
  
\begin{prob}
  $\Pow{\{a, b, c\}}$ అనే సమితిపై $\subseteq$ సంబంధంలోని
  \tecase{acc}{!!{element}s} జాబితాగా రాయండి.
\end{prob}

\end{document}
